\documentclass[12pt]{article}
\usepackage{sbc-template,graphicx,booktabs}
\usepackage{iftex}
\ifXeTeX
\IfFileExists{/usr/share/fonts/opentype/urw-base35/NimbusRoman-Regular.otf}{
\usepackage{fontspec}
\setmainfont[Path=/usr/share/fonts/opentype/urw-base35/,UprightFont=*-Regular,BoldFont=*-Bold,ItalicFont=*-Italic,BoldItalicFont=*-BoldItalic,Extension=.otf]{NimbusRoman}
\setsansfont[Path=/usr/share/fonts/opentype/urw-base35/,UprightFont=*-Regular,BoldFont=*-Bold,ItalicFont=*-Italic,BoldItalicFont=*-BoldItalic,Extension=.otf]{NimbusSans}
}{}
\fi
\usepackage[hidelinks]{hyperref}
\renewcommand{\tablename}{Tabela}
\renewcommand{\figurename}{Figura}
\newcommand{\tabela}[2]{\begin{center}\begin{minipage}{\linewidth}\captionof{table}{#1}\centering\small #2\end{minipage}\end{center}}
\newcommand{\figura}[3]{\begin{center}\begin{minipage}{\linewidth}\centering\includegraphics[width=#1\linewidth]{#2}\captionof{figure}{#3}\end{minipage}\end{center}}
\title{Classificação dos resultados do Santa Cruz de 2015 a 2025 com informações anteriores às partidas}
\author{Gabriel Renan Silva Martins Guimarães,\\ Anthonny Lucas Santos da Silva, Eliel dos Santos Rocha}
\address{Unidade Acadêmica de Belo Jardim -- Universidade Federal Rural de Pernambuco\\Belo Jardim -- PE -- Brasil\\Inteligência Artificial -- Prof. Dr. Anderson Pinheiro Cavalcanti}
\begin{document}
\maketitle
\begin{abstract}
We compared six classifiers on 468 Santa Cruz football matches using eight pre-match features. We tested three settings per classifier through stratified ten-fold cross-validation. Logistic regression achieved the highest mean macro F1 (0.420), whereas a shallow decision tree achieved the highest accuracy (0.453) without correctly identifying a single draw. Venue carried useful information, but identifying draws remained difficult across matches from different seasons.
\end{abstract}
\begin{resumo}
Comparamos seis classificadores em 468 partidas do Santa Cruz a partir de oito atributos disponíveis antes do jogo. Avaliamos três configurações por algoritmo em validação cruzada estratificada com dez partições. A regressão logística alcançou o maior F1 macro médio (0,420). Uma árvore rasa obteve a maior acurácia (0,453), mas errou todos os empates. O mando de campo ajudou a distinguir vitórias de derrotas. A identificação dos empates permaneceu uma dificuldade mesmo com os oito atributos.
\end{resumo}
\section{Introdução}
Antes de uma partida, sabe-se onde o Santa Cruz jogará, quem será o adversário e em que competição ocorrerá o encontro. Uma pergunta para a aprendizagem de máquina é quanto essas informações permitem distinguir três resultados do ponto de vista do clube: vitória, empate e derrota. Uma previsão de vitória fora de casa, por exemplo, representa uma vitória do visitante. Manter o ponto de vista do Santa Cruz é essencial desde a leitura da planilha até a interpretação dos erros.

O objetivo deste trabalho é comparar seis algoritmos para essa tarefa. Investigamos qual configuração reconhece melhor o conjunto das três classes, quanto o desempenho varia entre grupos de partidas e que diferença o mando faz quando as demais condições do experimento são mantidas. A análise dá continuidade à primeira entrega, na qual o grupo organizou a base e examinou a distribuição dos resultados.
\clearpage
\section{Trabalhos relacionados}
Bunker e Susnjak (2022) revisaram estudos sobre resultados em esportes coletivos. Os autores destacam a importância da escolha dos atributos e a dificuldade de comparar acurácias entre estudos com esportes, bases e condições diferentes. Neste projeto, dois cuidados vêm dessa observação. O primeiro é definir as informações disponíveis no momento do palpite. O segundo é comparar todas as configurações exatamente nos mesmos jogos de teste.

Berrar, Lopes e Dubitzky (2024) estudaram partidas futuras de várias equipes no \textit{Soccer Prediction Challenge}. O trabalho considerou tanto o placar exato quanto as probabilidades dos três resultados e explorou a evolução das equipes ao longo do tempo. KNN e redes neurais estiveram entre os métodos com bom desempenho naquele desafio. Aqui o recorte é um único clube e a saída é uma classe. A comparação que faz sentido é verificar se esses métodos também se destacam dentro da nossa base, com os atributos disponíveis nela.

Bunker, Yeung e Fujii (2024) discutem bases, atributos e formas de avaliação específicas do futebol. Essa distinção ajuda a situar o alcance deste estudo. Uma data ou uma fase de campeonato descreve o contexto de um jogo, mas informa pouco sobre a força das equipes naquele dia. Por isso, além de comparar algoritmos, examinamos os erros por classe e a frequência com que cada adversário aparece.
\section{Conjunto de dados}
A base \textit{Dados do time Santa Cruz Futebol Clube (2015--2025)} é disponibilizada no Kaggle (Guimarães, s.d.). Usamos a planilha com abas por temporada e conferimos os registros com o PDF organizado da primeira entrega. A versão tabular utilizada aqui reúne 470 partidas, com datas distintas, e 15 colunas. Ano e linha de origem são inteiros. As demais colunas guardam texto, identificadores, datas e horários conforme a organização da planilha.

A fonte tem duas maneiras de indicar as equipes. Algumas tabelas trazem mando e adversário. Outras identificam diretamente o time da casa e o visitante. Por exemplo, Macaé 2 a 0 Santa Cruz em 9 de maio de 2015 representa uma derrota do Santa Cruz fora de casa. Já um 2 a 0 com o Santa Cruz como mandante representa vitória. A inversão dessa lógica produziria rótulos errados em parte da amostra.

Na tabela organizada, 65 jogos não traziam o resultado na coluna destinada à classe, embora houvesse placar. Dez linhas tinham os conteúdos de competição e placar trocados. A correção foi feita antes de calcular o saldo de gols. Nos casos com disputa por pênaltis, o rótulo corresponde ao placar anterior às cobranças. Assim, um jogo empatado que termina em classificação nos pênaltis continua sendo um empate para a tarefa.
\clearpage
\subsection{Amostra e atributos}
Dois registros apresentavam divergência entre o resultado informado e o placar, fora de disputas por pênaltis: as partidas de 18 e 25 de junho de 2016. Sem uma confirmação externa, mantivemos o critério da primeira entrega e excluímos ambas. Restaram 468 exemplos em três classes: 181 vitórias (38,68\%), 143 empates (30,56\%) e 144 derrotas (30,77\%). A classe majoritária é a vitória. Em número de jogos, as duas menores classes diferem por apenas uma partida.

\tabela{Entradas dos modelos após a limpeza. A data completa foi mantida para auditoria.}{
\begin{tabular}{lll r}
\toprule Atributo & Representação & Valores observados & Ausentes\\\midrule
Ano & Inteiro & 2015 a 2025 & 0\\
Mês & Inteiro & 1 a 12 & 0\\
Dia da semana & Inteiro & 0 (segunda) a 6 (domingo) & 0\\
Horário & Minutos desde 0h & 660 a 1.365 & 11\\
Mando & Categoria & Casa ou fora & 0\\
Adversário & Categoria & 113 nomes & 0\\
Competição & Categoria & 40 nomes & 11\\
Fase & Categoria & 61 nomes & 4\\
\bottomrule\end{tabular}}

O número de partidas varia entre as temporadas. Em 2016 há 70 registros válidos, enquanto em 2024 há 13. Entre os adversários, o Náutico aparece 33 vezes, o Sport 30 e o Salgueiro 21. No outro extremo, 61 dos 113 adversários aparecem no máximo duas vezes. Para um jogo contra um desses times, o modelo pode ter apenas um exemplo do confronto no treino, ou nenhum.

\figura{0.72}{../resultados/mando.pdf}{Resultados por mando. Dos 233 jogos em casa, 118 terminaram em vitória. Fora de casa, foram 63 vitórias em 235 jogos.}
\clearpage
\section{Metodologia experimental}
\subsection{Preparação das entradas}
Usamos quatro atributos numéricos e quatro categóricos (Tabela 1). Para os numéricos, preenchemos ausências com a mediana do treino e aplicamos a transformação mínimo--máximo. Isso evita que diferenças de centenas de minutos dominem, pela magnitude, uma distância entre dois jogos. Ano, mês e dia da semana mantêm sua codificação numérica da primeira entrega. Ela simplifica a análise, mas deixa dezembro e janeiro distantes na escala, apesar de serem meses consecutivos.

Para as categorias ausentes, usamos o valor “Ausente”. O \textit{one-hot encoding} transforma cada categoria observada no treino em uma coluna indicadora. Assim, o algoritmo pode distinguir Sport de Náutico sem atribuir uma ordem numérica aos clubes. Se um adversário aparece pela primeira vez no teste, suas colunas de adversário ficam zeradas. O jogo ainda pode ser classificado com os demais atributos. Removemos espaços nas extremidades dos textos. Mantivemos as designações de fase e competição da fonte, inclusive quando a temporada faz parte do nome.

A data completa serve para identificar os jogos nos arquivos de resultados. Placar, gols e resultado informado servem para a auditoria da classe. Se o saldo de gols entrasse no modelo, uma simples comparação com zero resolveria a tarefa depois de encerrada a partida. Por isso, o treinamento recebe apenas as oito entradas da Tabela 1.
\subsection{Algoritmos e parâmetros}
Utilizamos implementações do scikit-learn (Pedregosa et al., 2011). A descrição dos modelos segue James et al. (2023). Para cada família, variamos um parâmetro em três configurações e mantivemos os demais fixos. Isso permite relacionar a mudança observada ao parâmetro que foi alterado.

\textbf{Árvore de Decisão.} A árvore divide os jogos por condições sobre os atributos e atribui uma classe a cada folha. Com o critério de Gini, a escolha de uma divisão favorece grupos com maior concentração de uma classe. Testamos profundidades máximas 2, 4 e 8. Um limite baixo produz regras mais gerais.

\textbf{Vizinhos Mais Próximos (KNN).} A previsão é a classe mais frequente entre os jogos de treino mais próximos. Utilizamos distância euclidiana e voto uniforme. Comparamos 3, 9 e 21 vizinhos. Aumentar a vizinhança reduz a influência de um único confronto, mas pode misturar contextos esportivos diferentes.

\textbf{Naive Bayes gaussiano.} O modelo estima a distribuição de cada entrada em cada classe e combina as evidências sob uma hipótese de independência condicional. Testamos suavização de variância de 0,01, 0,1 e 1. Esse ajuste evita que atributos com variância muito pequena exerçam peso excessivo. Uma distribuição gaussiana é uma aproximação limitada para indicadores que só assumem zero ou um.
\clearpage
\textbf{Regressão Logística.} Na versão multiclasse, o algoritmo combina as entradas com pesos aprendidos e atribui uma probabilidade a cada resultado. A previsão é a classe com maior valor. Usamos penalização L2 e testamos C igual a 0,1, 1 e 10. Valores menores impõem maior regularização aos pesos, o que pode ter importância para adversários com poucas observações.

\textbf{Máquina de Vetores de Suporte (SVM).} O algoritmo busca uma separação entre classes com margem e admite erros no treino. Usamos o núcleo radial (RBF), com o ajuste padrão de escala do parâmetro gamma. Variamos C entre 0,1, 1 e 10. Um C baixo permite maior tolerância a erros. O scikit-learn treina comparações entre pares de classes para o problema multiclasse.

\textbf{Rede Neural MLP.} A rede aprende combinações dos atributos em camadas. Testamos uma camada oculta de 8 neurônios, uma de 16 e duas de 16 e 8. Usamos ativação ReLU nas camadas ocultas, otimizador Adam, penalização L2 com alpha igual a 1, tolerância de 0,0005 e limite de 1.800 épocas. As três arquiteturas receberam as mesmas entradas.
\subsection{Validação e medidas}
Ordenamos os jogos por data e criamos dez partições estratificadas, com embaralhamento e semente 42. A estratificação mantém em cada parte proporções próximas de vitórias, empates e derrotas. Em cada rodada, nove partes servem para treino e uma para teste. O processo termina quando cada partida recebeu exatamente uma previsão fora do treino por configuração. São oito testes com 47 jogos e dois com 46, totalizando 180 ajustes das 18 configurações. Repetimos o procedimento com um classificador de referência que sempre escolhe a classe mais frequente do treino.

A mediana, os limites da escala e o vocabulário das categorias são calculados novamente a partir das nove partes de treino. Isso impede que, por exemplo, um horário do teste altere a escala usada no ajuste. Avaliamos a grade completa com um laço explícito no notebook, preservando as previsões e as métricas de cada rodada.

A acurácia é a fração de jogos classificados corretamente. A precisão de uma classe responde quantos palpites daquela classe estavam certos. A revocação (\textit{recall}) responde quantos resultados reais daquela classe o modelo recuperou. O F1 é a média harmônica entre precisão e revocação. Quando o denominador é zero, atribuímos zero à medida. Calculamos precisão, revocação e F1 para cada classe e tiramos a média simples das três (macro). Assim, o empate tem o mesmo peso que a vitória.

Para cada configuração, reportamos a média aritmética das dez medidas e o desvio padrão amostral, com divisor nove. Adotamos o F1 macro médio para comparar configurações. A escolha da maior média é exploratória: os mesmos folds participam da comparação dos parâmetros, e o desempenho da configuração escolhida ainda precisa de confirmação em uma amostra externa.
\clearpage
\section{Resultados}
A Tabela 2 apresenta as 18 configurações na mesma escala. Cada célula contém a média em porcentagem e o desvio padrão em pontos percentuais, calculados sobre dez testes. Na linha geral, primeiro calculamos, em cada fold, a média das 18 configurações. Depois resumimos esses dez valores. A referência majoritária é mostrada à parte e fica fora da média geral do dataset.

\tabela{Desempenho no dataset Santa Cruz. Precisão, recall e F1 são médias macro. O ajuste é o parâmetro variado em cada família, descrito na metodologia.}{
\setlength{\tabcolsep}{4pt}\begin{tabular}{llrrrr}\toprule
Modelo & Ajuste & Acurácia & Precisão & Recall & F1\\\midrule
Árvore & 2 & 45,3 ± 5,9 & 30,2 ± 4,0 & 43,3 ± 5,7 & 35,3 ± 4,5\\
 & 4 & 42,7 ± 4,8 & 31,3 ± 6,9 & 41,0 ± 5,0 & 34,5 ± 4,2\\
 & 8 & 44,0 ± 7,1 & 40,4 ± 9,5 & 42,8 ± 7,3 & 39,2 ± 7,3\\
\midrule
KNN & 3 & 40,0 ± 9,1 & 38,8 ± 10,6 & 39,7 ± 9,3 & 37,9 ± 10,1\\
 & 9 & 41,2 ± 10,0 & 39,2 ± 11,1 & 40,1 ± 10,3 & 39,1 ± 10,5\\
 & 21 & 43,6 ± 6,6 & 39,9 ± 8,6 & 41,9 ± 6,8 & 39,8 ± 7,1\\
\midrule
Naive Bayes & 0,01 & 32,7 ± 5,3 & 34,1 ± 6,3 & 33,4 ± 5,5 & 32,7 ± 5,4\\
 & 0,1 & 38,5 ± 7,0 & 39,5 ± 7,2 & 38,9 ± 6,9 & 38,4 ± 6,9\\
 & 1,0 & 42,5 ± 5,5 & 39,6 ± 6,1 & 40,6 ± 5,3 & 38,8 ± 5,0\\
\midrule
Logística & 0,1 & 42,1 ± 5,7 & 36,2 ± 5,5 & 39,8 ± 5,3 & 36,8 ± 4,8\\
 & 1 & 44,3 ± 6,6 & 42,8 ± 6,8 & 42,7 ± 6,5 & 42,0 ± 6,6\\
 & 10 & 40,8 ± 9,9 & 40,2 ± 10,1 & 40,1 ± 9,8 & 39,6 ± 9,8\\
\midrule
SVM & 0,1 & 38,7 ± 0,7 & 12,9 ± 0,2 & 33,3 ± 0,0 & 18,6 ± 0,2\\
 & 1 & 42,3 ± 5,2 & 38,4 ± 6,7 & 40,1 ± 5,4 & 38,2 ± 5,8\\
 & 10 & 38,3 ± 9,4 & 37,2 ± 9,5 & 37,5 ± 9,6 & 36,9 ± 9,5\\
\midrule
MLP & 8 & 43,8 ± 6,7 & 39,9 ± 7,1 & 41,9 ± 6,8 & 39,9 ± 7,1\\
 & 16 & 44,0 ± 6,2 & 40,8 ± 8,0 & 42,2 ± 6,0 & 40,5 ± 6,7\\
 & 16 e 8 & 39,8 ± 6,4 & 35,7 ± 10,4 & 38,2 ± 6,6 & 36,1 ± 8,9\\
\midrule
Majoritária & constante & 38,7 ± 0,7 & 12,9 ± 0,2 & 33,3 ± 0,0 & 18,6 ± 0,2\\
\midrule Geral & 18 ajustes & 41,4 ± 4,9 & 36,5 ± 5,4 & 39,9 ± 5,0 & 36,9 ± 5,1\\\bottomrule\end{tabular}}

A referência obteve 38,7\% de acurácia. Como todas as previsões foram de vitória, a revocação dessa classe foi de 100\% e a das outras duas foi zero. A média macro das três revocações ficou em 33,3\%. A SVM com C igual a 0,1 repetiu exatamente as previsões da referência nos 468 jogos. Com C igual a 1, a SVM passou a identificar exemplos das três classes, e seu F1 macro subiu de 18,6\% para 38,2\%.

A árvore de profundidade 2 atingiu a maior acurácia média, de 45,3\%. Das 468 previsões, apenas uma foi de empate, e ela estava errada. Ao permitir oito níveis, a árvore acertou 20 empates, mas perdeu parte dos acertos em vitórias e derrotas. Essa mudança elevou o F1 macro mesmo com a queda da acurácia. A regressão logística com C igual a 1 alcançou o maior F1 macro (42,0\%) e a maior precisão macro (42,8\%).
\clearpage
\subsection{Variação entre modelos e erros por classe}
A Figura 2 reúne a configuração de maior F1 macro de cada família. A regressão logística ocupa a primeira posição, seguida da MLP com 16 neurônios. A diferença de aproximadamente 1,5 ponto percentual entre as duas médias é pequena em relação à variação de cada modelo entre folds. As barras descrevem a dispersão observada. Como as amostras de treino se sobrepõem, a figura não permite concluir que a diferença seja estatisticamente significativa.

\figura{0.85}{../resultados/comparacao.pdf}{F1 macro das configurações selecionadas e da referência majoritária. O ponto indica a média e a barra representa um desvio padrão amostral.}

\tabela{Matriz de confusão da regressão logística com C igual a 1. As linhas indicam o resultado real e as colunas, a previsão. Cada partida é contada uma vez.}{
\begin{tabular}{lrrrr}\toprule
Real / previsto & Vitória & Empate & Derrota & Total\\\midrule
Vitória & 110 & 44 & 27 & 181\\
Empate & 63 & 34 & 46 & 143\\
Derrota & 41 & 40 & 63 & 144\\\midrule
Total & 214 & 118 & 136 & 468\\\bottomrule
\end{tabular}}

Dos 143 empates, apenas 34 foram reconhecidos (23,8\%). A maioria dos erros nessa linha foi classificada como vitória. As taxas de reconhecimento agregadas de vitórias e derrotas foram 60,8\% e 43,8\%, respectivamente. O modelo fez 118 previsões de empate e acertou 34 delas, o que corresponde a 28,8\% de precisão para a classe. Portanto, a dificuldade envolve tanto encontrar os empates reais quanto evitar alarmes falsos.

A proximidade entre as quantidades de empates e derrotas ajuda a interpretar essa diferença. Houve 143 exemplos de uma classe e 144 da outra, mas o modelo reconheceu quase o dobro de derrotas. A frequência das classes, sozinha, não explica o resultado. Os atributos disponíveis separam melhor algumas situações de vitória e derrota do que as situações que terminam em igualdade.
\clearpage
\section{Discussão}
\subsection{Quanto o mando acrescenta?}
Em casa, o Santa Cruz venceu 50,6\% das partidas. Fora, venceu 26,8\%. Para investigar o peso dessa diferença, repetimos a regressão logística com C igual a 1 nos mesmos folds, mudando somente o conjunto de atributos. Mantivemos o pré-processamento correspondente às entradas de cada teste. A Tabela 4 mostra o que acontece ao retirar o mando e ao utilizá-lo isoladamente.

\tabela{Efeito do mando na regressão logística. Média em porcentagem e desvio padrão em pontos percentuais, calculados nos mesmos dez folds.}{
\begin{tabular}{lcc}\toprule Entradas & Acurácia & F1 macro\\\midrule
Somente mando & 44,9 ± 6,3 & 35,0 ± 4,9\\
Sete atributos sem o mando & 40,4 ± 6,6 & 39,2 ± 7,1\\
Oito atributos, incluindo mando & 44,3 ± 6,6 & 42,0 ± 6,6\\\bottomrule\end{tabular}}

Usar somente o mando favorece a acurácia porque o modelo tende a escolher vitória em casa e derrota fora. Essa regra deixa o empate sem uma previsão própria. Ao incluir os demais atributos, o F1 macro subiu de 35,0\% para 42,0\%. No sentido inverso, retirar apenas o mando da base completa reduziu o F1 macro em 2,8 pontos percentuais. Essa comparação mede a contribuição do atributo para este classificador e esta amostra. Estimar um efeito causal de jogar em casa exigiria outro desenho de pesquisa.

\subsection{Parâmetros, representação e limites}
Com três vizinhos, um único jogo representa um terço dos votos do KNN. Com 21, a votação depende menos de uma observação isolada. O aumento de 37,9\% para 39,8\% no F1 macro é compatível com essa interpretação. Após a codificação, dois adversários ocupam posições distintas, mas os indicadores não expressam se as equipes têm forças parecidas.

Para testar diretamente o efeito da escala, repetimos o KNN com 21 vizinhos sem a normalização numérica. A acurácia caiu de 43,6\% para 33,3\%, e o F1 macro, de 39,8\% para 30,8\%. Com as demais condições fixas, a escala mudou quais partidas o algoritmo considerou semelhantes.

No Naive Bayes, aumentar a suavização elevou o F1 macro de 32,7\% para 38,8\%. Um indicador de adversário raro pode ter variância próxima de zero em uma classe. A suavização limita seu peso. Na regressão logística, a melhor média ocorreu com regularização intermediária. Já a MLP de duas camadas teve desempenho inferior às alternativas de uma camada. Mais parâmetros não trouxeram ganho nesta amostra. 

A validação estratificada atende ao protocolo da disciplina e usa todos os jogos para teste. Como as temporadas foram misturadas, um jogo de 2017 pode ser classificado por um modelo que treinou com jogos de 2025. Os resultados representam a classificação de partidas retiradas desse histórico. Para estudar a previsão da próxima temporada, seria necessário manter o treino no passado e reservar jogos posteriores para o teste.
\clearpage
\section{Considerações finais}
Entre as 18 configurações, a regressão logística com C igual a 1 apresentou o maior F1 macro médio. Se o objetivo fosse apenas maximizar o total de palpites corretos, a árvore de profundidade 2 lideraria a comparação. Seus erros mostram a consequência dessa escolha: nenhum dos 143 empates foi reconhecido. Para uma tarefa que atribui importância às três classes, a matriz de confusão e o F1 macro mudam a conclusão que a acurácia isolada sugeriria.

O mando ajuda a classificar os jogos, mas possui poder limitado para antecipar empates. O calendário, o nome dos clubes e a fase da competição também descrevem pouco do desempenho recente. Uma continuação viável seria calcular, para cada equipe, pontos e saldo de gols em partidas anteriores, usando apenas dados já disponíveis na data do palpite. Esses atributos poderiam ser comparados aos atuais em um teste cronológico com os últimos anos reservados. Para separar a escolha dos parâmetros da avaliação final, a seleção poderia ocorrer exclusivamente nas temporadas de treino.

O material que acompanha o artigo contém a planilha original, o PDF organizado da primeira entrega, as tabelas antes e depois da limpeza, o código-fonte e os resultados por fold. A tabela de previsões permite conferir uma partida específica pela data e reconstruir a matriz de confusão. O notebook também reproduz a média geral do dataset e a comparação dos três conjuntos de atributos utilizados na análise do mando.
\section*{Referências}
\noindent Berrar, D., Lopes, P. e Dubitzky, W. (2024). A data- and knowledge-driven framework for developing machine learning models to predict soccer match outcomes. \textit{Machine Learning}, 113, 8165--8204. DOI: 10.1007/s10994-024-06625-9.

\noindent Bunker, R. e Susnjak, T. (2022). The application of machine learning techniques for predicting match results in team sport: a review. \textit{Journal of Artificial Intelligence Research}, 73. DOI: 10.1613/jair.1.13509.

\noindent Bunker, R., Yeung, C. e Fujii, K. (2024). Machine learning for soccer match result prediction. Em \textit{Artificial Intelligence, Optimization, and Data Sciences in Sports}, p. 7--49. Springer. DOI: 10.1007/978-3-031-76047-1\_2.

\noindent Guimarães, G. R. S. M. (s.d.). \textit{Dados do time Santa Cruz Futebol Clube (2015--2025)}. \href{https://www.kaggle.com/datasets/gabrielrenanpe/dados-do-time-santa-cruz-futebol-clube-2015-2025}{Kaggle, perfil gabrielrenanpe}.

\noindent James, G., Witten, D., Hastie, T., Tibshirani, R. e Taylor, J. (2023). \textit{An Introduction to Statistical Learning: with Applications in Python}. Springer. DOI: 10.1007/978-3-031-38747-0.

\noindent Pedregosa, F. et al. (2011). Scikit-learn: machine learning in Python. \textit{Journal of Machine Learning Research}, 12, 2825--2830.
\end{document}
